\begin{abstract}
Instruction-based image editors such as InstructPix2Pix let users apply a sequence of natural-language edits to an image, yet no standard metric captures what happens to the parts of the image nobody asked to change. This paper introduces the Drift Score, a metric designed to be model-agnostic that combines Segment Anything (SAM) region segmentation with CLIP embedding similarity to quantify unintended change after an edit. The metric is applied across 120 manually constructed four-step edit chains built from 60 COCO images, split evenly between localized object-level edits and broad scene-level edits. Two low-cost mitigation strategies, masked conditioning and region-locking, are implemented and evaluated against an unmitigated baseline using paired significance testing, alongside an independent Edit Adherence Score that checks whether the requested edit still happened. Both mitigations produce large, highly significant reductions in measured drift (region-locking: 89.1\%, $p<0.0001$, $d_z=2.68$; masked conditioning: 68.6\%, $p<0.0001$, $d_z=1.73$), at a comparatively smaller but still significant cost to instruction adherence. Two hypotheses formulated before the experiment did not hold as predicted: region-locking significantly outperforms masked conditioning rather than the reverse, and the data shows no reliable evidence of drift increasing across the chain, with the only significant step-position effect running opposite to the expected compounding pattern. Both reversals are traced to specific mechanisms rather than left unexplained. The Drift Score, dataset, and full pipeline are released as a reproducible toolkit built entirely from pretrained models.
\end{abstract}
\begin{IEEEkeywords}
Image Editing, Diffusion Models, Semantic Drift, CLIP, Segment Anything, Evaluation Metrics, Instruction-Based Editing
\end{IEEEkeywords}
